% ! TeX root = ./main.tex

\section{Lessons From Our Placement Exercise}
\label{sec:amr:lessons}

% Mirroring Petrini et al.'s~\cite{Petri2003:Missing} findings on flagship systems two decades ago,
% our investigation underscores the persistent need for deep, cross-stack analysis to diagnose
% hidden AMR performance bottlenecks. 
We distill our takeaways from this exercise into the following five lessons.

\lesson{1}{AMR Performance Is Straggler Mitigation}

High synchronization costs in AMR codes are often attributed to load imbalance.
Our experience suggests that this attribution may be premature: stragglers can also arise from
tuning artifacts, execution strategy interactions, or hardware-level variability.
Placement strategies cannot compensate for unstable system behavior---such instability must be
identified and resolved independently before placement can be meaningfully evaluated.

\lesson{2}{Effective Tuning is Empirical and Idiosyncratic}

Performance tuning is not about correctness, but about aligning system behavior with the workload.
Default configurations reflect common-case assumptions that may vary across codes.
Tuning, therefore, is both empirical and
workload-specific~\cite{Kruse2023:StatisticalAnalysisHPC,Mijak2016:PERISCOPE,Pelle2012:TuningMPI}.
It improves performance directly, while also making application behavior more predictable under
placement changes.
In our experience, progressive tuning clarified the telemetry structure, reduced confounding noise,
and exposed deeper residual issues that would otherwise be misattributed.

\lesson{3}{P2P Communication Is A Key Variability Source}

Our tuning efforts primarily focused on point-to-point (P2P) communication patterns used in ghost
cell and halo exchanges.
Unlike collectives, which are structured and easy to benchmark in isolation, P2P exchanges are
fine-grained, latency-sensitive, and highly variable.
They involve dynamic neighbor counts, varying message sizes, and inconsistent communication
paths---each of which can activate different code paths in the network stack.
We found these interactions to be the dominant source of variability, and they required targeted
attention during tuning.
% While newer MPI standards provide neighbor collectives to optimize such patterns, these mechanisms
% often fall short of fully stabilizing P2P behavior in AMR codes.

\lesson{4}{Diagnosis Needs Structured, Queryable Telemetry}

Understanding why performance deviated required analytics over fine-grained telemetry to surface
anomalous low-level behavior.
% citation for JSON (chrome trace)
While tracing tools capture such detail, they rely on unstructured formats like
JSON~\cite{Googl:TraceEventFmt} or OTF2~\cite{Eschw2012:OTF2}, which are poorly suited to
multi-dimensional analysis across rank, time, and task.
We converged on SQL-based analytics over telemetry grouped by timestep and sorted by rank, enabling
views of application behavior aligned with synchronization intervals.

While tools like Caliper~\cite{Boehm2016:Caliper} do support structured telemetry, they---like our
ad hoc pipeline---rely on parsing and postprocessing of plaintext data, which we found to slow down
our iterative tuning workflows.
Binary columnar formats like Arrow~\cite{:ApacheArrow} and Parquet~\cite{:ApacheParquet}, when
paired with in-situ collection, offer a promising foundation for low-latency BSP telemetry by
enabling low-overhead parsing and efficient querying via embedded statistics over partitioned data.

\lesson{5}{Placement Optimization is Systems Engineering}

% Designing placement policies in AMR codes is a systems engineering problem, not a theoretical
% exercise. AMR placement is a systems engineering problem grounded in empirical constraints rather
% than algorithmic considerations.
AMR placement policy design is a systems engineering problem, guided primarily by empirical
constraints rather than solely by theoretical optimality.
Tradeoffs between compute balance, communication locality, and placement overhead must be evaluated
based on observed performance impact, not static preferences.
% Universal strategies or fixed priorities are rarely effective; each deployment demands empirical
% calibration.
Placement strategies need to be empirically calibrated to each specific deployment context.
Their computation must also respect strict time budgets---in our case, under 50\,ms per
redistribution---which constrains the complexity of the algorithms and necessitates focus on
dominant effects.
